\documentclass[11pt]{article}

\usepackage[utf8]{inputenc}
\usepackage[T1]{fontenc}
\usepackage{lmodern}
\usepackage[margin=1in]{geometry}
\usepackage{amsmath,amssymb}
\usepackage{graphicx}
\usepackage{hyperref}
\usepackage{booktabs}
\usepackage{microtype}
\usepackage{xcolor}
\usepackage{listings}
\usepackage{titlesec}
\usepackage{caption}
\usepackage{enumitem}
\usepackage{authblk}
\usepackage{url}

\hypersetup{
  colorlinks=true,
  linkcolor=blue!60!black,
  citecolor=blue!60!black,
  urlcolor=blue!60!black,
  pdftitle={The Thalamus Layer: Native Vector Routing for Edge AI Agents},
  pdfauthor={Muhammet Sıddık Bel}
}

\titleformat{\section}{\normalfont\large\bfseries}{\thesection}{0.6em}{}
\titleformat{\subsection}{\normalfont\normalsize\bfseries}{\thesubsection}{0.6em}{}

\title{\Large\bfseries The Thalamus Layer:\\
Native Vector Routing for Edge AI Agents}

\author[1]{Muhammet Sıddık Bel\thanks{\texttt{msbel5@gmail.com}}}
\affil[1]{Independent. Diyarbakır, Turkey. \url{https://msbel.com}}

\date{April 29, 2026}

\begin{document}
\maketitle

\begin{abstract}
\noindent
Multi-agent large language model systems serialize all inter-module
communication through human language. Vision encoders produce dense
embeddings, then caption them into English. Planning modules read the
caption, re-embed it, plan, and write a sentence to the critic. The
critic re-embeds and produces a verdict. By the time a single decision
is made, the same content has been encoded into and decoded out of
embedding space five or six times. We argue that this is the dominant
hidden cost of contemporary agentic systems. We propose an alternative
runtime architecture in which specialist encoders communicate via
modality-aligned vectors through a software router we call the
\emph{thalamus layer}, with a layered hippocampal memory that preserves
modality and an idle replay loop modeled on the default mode network.
We describe a Phase 1 reference implementation as an OpenClaw plugin,
deployable on a Raspberry Pi 5 with 4 GB of RAM, and we sketch a
two-step roadmap toward distributed and hardware-accelerated versions.
We position this proposal against contemporary work on latent
multi-agent collaboration, perceiver-style architectures, and
neurobiologically inspired retrieval, and argue that the missing
contribution is not a new model but a small, open, replicable runtime
that runs on hardware a hobbyist can own.
\end{abstract}

\noindent\textbf{Keywords:} multi-agent systems, latent space,
LLM agents, OpenClaw, edge inference, agent runtime, hippocampal
memory, default mode network.

\section{Introduction}

The unit of communication between modules in modern multi-agent
language model systems is text. A vision encoder produces a
high-dimensional embedding that captures texture, geometry, and
material cues. The system then collapses this embedding into a
short English caption such as ``a corroded pipe with visible
water leak.'' A downstream planning model re-embeds the caption,
loses most of the original visual fidelity, and proceeds. A
critic reads the planner's text output, re-embeds again, and
produces a verdict. None of these encode-decode round trips are
required for the consumers to function. They are an artifact of
an interface convention inherited from the chatbot era, in which
the only consumer was a human and the only viable medium was
language.

The cost of this convention is mostly invisible because it
distributes itself across many small budget items: tokens spent
generating intermediate captions, latency added by extra forward
passes, fidelity discarded at each modality boundary. Aggregated,
these costs dominate latency and per-task expense in non-trivial
agent loops. They also place a hard ceiling on what small models
can do in agent settings, because each round trip leaves less
budget for actual reasoning.

This paper argues that the next architectural step for agentic
systems is to remove text from the inter-module bus when the
consumer is another module, and reserve text for the final
user-facing turn. We refer to the routing layer required to make
this work as the \emph{thalamus layer}, by analogy to the
biological thalamus, which serves as a relay and gating structure
between sensory cortex and cortical processing areas without
itself performing reasoning. We describe a memory layer modeled
on the hippocampus, an idle-time replay loop modeled on the
default mode network (DMN), and a roadmap toward distributed and
hardware-accelerated versions of the same architecture.

The contribution of this paper is not a new model. It is an
\emph{architectural proposal supported by a working Phase 1
implementation}, deployable on a single Raspberry Pi 5, released
under the MIT license, with companion plugins for tool-call audit
and atom retrieval already published as installable npm packages.

\section{The Bottleneck}
\label{sec:bottleneck}

Consider a typical agentic flow that processes a multimodal user
query, for instance a photo of a broken appliance accompanied by
a request for repair instructions.

\begin{enumerate}
  \item The vision encoder ingests the image and produces a
        $d_v$-dimensional embedding $v \in \mathbb{R}^{d_v}$,
        with $d_v$ in the range $512$ to $1024$ for current
        small encoders and up to $4096$ for larger ones.
  \item The system requires the embedding to flow downstream.
        Lacking a native vector channel, it asks the vision
        model to caption the image. The captioner emits
        $C_1 = \mathrm{decode}(v)$, a short English string.
  \item The planning model re-embeds $C_1$ to $v' \in
        \mathbb{R}^{d_p}$ where $d_p \neq d_v$ in general, and
        proceeds. The mapping $v \to C_1 \to v'$ is many-to-one
        followed by stochastic decode, and is therefore lossy by
        construction.
  \item The planning model emits a plan as text $C_2$, which the
        critic re-embeds, evaluates, and replies to as text $C_3$.
  \item The episode is committed to memory, often as $C_3$ alone,
        with the original embeddings $v, v'$ discarded.
\end{enumerate}

The information loss at each $v \to C \to v'$ boundary is not
hypothetical. Consider that $v$ may distinguish between the rust
patterns of cast iron and copper plumbing, while $C_1$ likely
flattens both to ``corroded pipe.'' The discriminative signal is
present in the encoder output and absent in the caption. A
downstream module that consumes only $C_1$ cannot recover it.

\subsection{A small empirical probe}

To make the bottleneck concrete, we generated three sentences
through OpenAI's \texttt{text-embedding-3-small} encoder
(1536 dimensions, $L_2$-normalized to unit norm) on a Raspberry
Pi 5, then presented the resulting vectors to Anthropic's Claude
Sonnet 4.6 along with the model name, dimensionality, and the
first 20 floats of each vector. We did not provide the original
sentences.

The model correctly identified the encoder family from the
dimension and normalization fingerprint, computed approximate
inter-sentence cosine alignment from the visible coordinates,
and reported the inputs as well-formed embeddings. It explicitly
stated that it could not recover the original sentences, that it
could not determine which embedding corresponded to which
sentence, and that the mapping from individual dimensions to
semantic content was opaque to it.

This is the bottleneck made visible. The information was present
in the float vectors, but the consumer was a model trained on
text and images, not on raw embeddings. The decoder pass that
would have recovered the sentence was not part of the consumer's
forward path. In a realistic agent flow, this is also why the
system cannot route the embedding directly: there is no
ingress for it on the consuming side.

\subsection{Cost analysis}

Let $\tau_e$ be the encoder forward pass cost, $\tau_d$ the
decoder forward pass cost, and $L$ the average length of an
intermediate text artifact in tokens. A naive encode-decode-
reencode hop incurs $\tau_d + L \cdot c_t + \tau_e$ where
$c_t$ is the per-token autoregressive cost. For a chain of
$k$ specialist modules under the current text-bus convention,
the aggregate inter-module cost is approximately
$(k-1)(\tau_d + L \cdot c_t + \tau_e)$, plus any text
context held by intermediate modules.

Under a native vector bus the equivalent cost is at most
$(k-1) \cdot \tau_a$, where $\tau_a$ is the cost of a small
adapter projection between modality spaces, typically a
single linear layer or shallow MLP with two to three orders of
magnitude fewer parameters than $\tau_d$ or $\tau_e$. The
difference is dominant in the typical regime where $L$ is
in the hundreds of tokens and $c_t$ is non-negligible.

The fidelity argument is independent. Even at $c_t \to 0$,
the chain remains lossy because each text serialization is a
many-to-one map.

\section{Architecture}
\label{sec:architecture}

We propose a runtime composed of six logical layers (Figure
\ref{fig:arch}).

\begin{figure}[ht]
\centering
\begin{verbatim}
                 user input
                     |
                     v
   +-------------------------------------+
   |  specialist encoders                |
   |  (vision, audio, text, ...)         |
   +-----------------+-------------------+
                     |  raw embeddings
                     v
   +-------------------------------------+
   |  modality adapters                  |
   |  (frozen or trained projections)    |
   +-----------------+-------------------+
                     |  shared workspace vectors
                     v
   +-------------------------------------+
   |  thalamus router                    |
   |  (gating, prioritization, FIFO)     |
   +--+-----+-----+-----------+----------+
      |     |     |           |
      v     v     v           v
   +--+--+ +-+--+ +--+----+ +-+-------+
   |reas | |crit| |planner| |memory   |
   |core | |core| |core   | |layer    |
   +-----+ +----+ +-------+ +---------+
                                |
                  hippocampal memory store
                                |
                  +-------------+-------------+
                  | hot (Redis) | semantic DB |
                  | episodic    | raw embed   |
                  +-------------+-------------+

                language head only at
                final user-facing turn
\end{verbatim}
\caption{The thalamus layer architecture. Specialist encoders
produce modality-specific embeddings, which are projected into
a shared workspace by frozen or trained adapters. The thalamus
router gates and forwards packets to consumer modules. A
layered hippocampal store preserves modality and supports
idle-time replay (DMN). Text is generated only at the final
user-facing step.}
\label{fig:arch}
\end{figure}

\subsection{Specialist encoders}
The specialist layer consists of unaltered, off-the-shelf
encoders for each modality the system supports. Vision uses
SigLIP \cite{siglip} or DINOv2 \cite{dinov2}. Audio uses a
Whisper encoder front-end \cite{whisper} or EnCodec
\cite{encodec}. Text uses a small embedding model such as
\texttt{e5-small} \cite{e5} or
\texttt{text-embedding-3-small}. None of these are modified.

\subsection{Modality adapters}
Each encoder produces vectors in a different ambient space and
of a different dimensionality. The adapter layer projects each
modality's output into a shared workspace vector space of fixed
dimension $d_w$ (default 512 in our reference implementation).
In Phase 1 these adapters are frozen random orthogonal linear
maps; the workspace is consistent but not aligned across
modalities. In Phase 1.1 they are trained on paired data
(image-caption, audio-transcript) with a contrastive loss to
maximize cross-modal cosine similarity.

\subsection{Thalamus router}
The router is a software loop that holds a queue of \emph{packets}.
Each packet has a unique identifier, a source modality, a source
module, an optional target module, the workspace vector, a
priority bucket, a creation timestamp, and a metadata dictionary.
The router applies FIFO ordering within each priority bucket,
drops low-priority packets when queue depth exceeds a threshold,
and enforces a maximum hop count of three to prevent cycles.
Routing decisions are logged in append-only JSONL for downstream
audit, in line with the AEGIS plugin pattern \cite{aegis}.

The router \emph{is not the smart part of the system}. It does
no reasoning. The biological analogy is intentional: the
thalamus relays signals to cortical regions and gates them by
attention, but cortical computation happens elsewhere. The
contribution of the router is to make the inter-module bus
modality-preserving and structurally observable.

\subsection{Reasoning, critic, and planner cores}
These are the cortical specialists. Each consumes workspace
vectors, produces workspace vectors, and may emit additional
metadata. In our reference implementation, reasoning and critic
cores are existing OpenClaw agents (a five-agent crew composed
of Liaison, Captain, Builder, Inspector, and Archivist roles)
exposed to the thalamus through small Python wrappers that
translate between agent message formats and packets.

\subsection{Hippocampal memory}
The memory layer is built from four components, each of which
already exists in the OpenClaw plugin ecosystem: a hot working-
memory tier in Redis, a vector index in LanceDB \cite{memlance},
an episodic markdown store in an Obsidian-compatible vault
\cite{memwiki}, and a fourth, modality-preserving raw embedding
cache. The fourth component is novel relative to the published
OpenClaw memory plugins and is the smallest concrete addition
required to operate the thalamus end-to-end.

\subsection{Default Mode Network replay}
The DMN component is an idle-loop replay mechanism that
activates when the agent has no active user-facing task. It
samples past packets from the embed cache, re-runs them through
the reasoning core, and writes new association rows when
results differ from the originally stored output. This is
functionally analogous to memory consolidation during slow-wave
sleep \cite{dmn-review}, and operationally it is a heartbeat
slot that takes at most one minute of wallclock per invocation.
It is bounded, observable, and terminates.

\section{Reference Implementation}
\label{sec:impl}

The reference implementation, named \texttt{openclaw-thalamus},
is a Phase 1 software simulation deployed as an OpenClaw plugin
on a Raspberry Pi 5 with 4 GB of RAM. It uses NumPy for vector
arithmetic and Torch (CPU only) for encoder forward passes. It
ships under MIT license at \url{https://github.com/msbel5/openclaw-thalamus}
and at \url{https://www.npmjs.com/package/@msbel/openclaw-thalamus}.

The plugin registers three MCP tools, \texttt{thalamus\_encode},
\texttt{thalamus\_route}, and \texttt{thalamus\_recall}, and
subscribes to the \texttt{agent\_end} hook to trigger DMN replay.
The plugin coexists with two earlier plugins from the same
author: \texttt{@msbel/openclaw-aegis-signer}, which provides
Ed25519-signed tool-call audit \cite{aegis}, and
\texttt{@msbel/openclaw-sga-mcts-atoms}, which provides
plan-time tool-use atom retrieval \cite{sgamcts}. Together
these three plugins constitute a small, observable, and
auditable agent runtime running entirely on commodity edge
hardware.

We emphasize that Phase 1 is single-process and does not yet
realize the full latency benefits of native vector routing. It
demonstrates the \emph{architecture} on hardware that is
accessible to individual researchers and hobbyists, and it
removes text from the inter-module bus for cases where the
consumer is another module.

\section{Related Work}
\label{sec:related}

\paragraph{Latent multi-agent collaboration.}
LatentMAS \cite{latentmas} proposes inter-agent collaboration in
latent space rather than text and reports favorable performance
in benchmark settings. Our proposal differs in placing latent
collaboration inside a single edge runtime with explicit
hippocampal memory and DMN replay, and in shipping a working
plugin rather than a benchmark.

\paragraph{Latent prediction and perceiver-style architectures.}
V-JEPA \cite{vjepa} performs video prediction directly in
representation space rather than at the pixel level, supporting
the design principle ``do not decode just to encode again.''
Perceiver IO \cite{perceiverio} is the architectural ancestor
for general-purpose latent processing across modalities and
informs the modality-adapter and shared-workspace decisions in
this paper.

\paragraph{Neurobiologically inspired retrieval.}
HippoRAG \cite{hipporag} models the hippocampal indexing
function via a knowledge graph plus Personalized PageRank for
multi-hop retrieval. Our hippocampal memory layer is more
modest, layering existing components, but the design intent is
aligned: the memory is an indexer between short-term and
long-term storage, not the central processor.

\paragraph{Global Workspace Theory in agents.}
Theater of Mind \cite{theaterofmind} proposes Global Workspace
Agents with a central broadcast hub plus heterogeneous
functionally-constrained agents, and is the closest existing
work to the thalamus and cortex split presented here. The
present paper differs in providing an edge implementation and
in committing to native vector exchange rather than broadcast
of structured messages.

\paragraph{Embodied and continuous cognition.}
JEPA \cite{jepa} and related work from the I-JEPA / V-JEPA
program articulate continuous prediction in representation
space as a route past the limits of token-level autoregression.
The thalamus layer is complementary, addressing inter-module
communication where the modules are not necessarily JEPA
predictors.

\section{Roadmap}
\label{sec:roadmap}

\paragraph{Phase 1: software simulation, in-process.} The
implementation described in Section~\ref{sec:impl}. Single
Python process, all encoders and adapters loaded as modules,
tensor exchange via NumPy in shared memory, memory layers as
already deployed, DMN loop as a heartbeat slot. Already
deployable on a 4 GB Pi 5.

\paragraph{Phase 2: distributed, multi-process.} Encoders,
reasoning cores, and memory each run as separate processes.
The thalamus runs in its own process and addresses peers via
shared memory or RDMA-style local transport. JSON metadata
only, with binary tensor payloads transmitted out of band. This
is where the latency improvements over the text bus become
empirically measurable. The hardware floor likely rises to an
8 GB Pi or a small cluster of Pi 5s.

\paragraph{Phase 3: hardware thalamus.} A dedicated routing
fabric, plausibly an FPGA card or an NPU operated as a routing
accelerator (the Hailo class of devices in the AI HAT+ family
is a candidate). The job of this hardware is to move tensors
between modules with predictable latency and gating semantics,
approximating the role of the biological thalamus. This phase
is more than a year out and requires collaborators.

\section{Discussion and Limitations}
\label{sec:discussion}

\paragraph{Adapter alignment.} The most important open
question in this proposal is whether modality adapters can be
made to produce a meaningfully shared workspace without
expensive paired-data training. Phase 1 punts on this question
by using random orthogonal projections, which guarantee
consistency but not alignment. Phase 1.1 must answer it with
contrastively trained adapters; if it fails, the architecture
collapses to ``faster but no smarter.''

\paragraph{Encoder choice constrains the shared workspace.}
SigLIP and CLIP-family encoders are pretrained against text and
have an implicit shared space already. Audio encoders such as
Whisper's encoder and EnCodec do not. The hardest adapters are
audio-to-workspace and reasoning-state-to-workspace.

\paragraph{The DMN replay loop is bounded.} Real biological DMN
runs continuously and at scale. The Phase 1 replay loop is a
heartbeat-bounded, one-minute-per-tick approximation. It
demonstrates the mechanism but does not approach the
qualitative behavior of biological consolidation. We do not
claim it does.

\paragraph{No claims about consciousness.} This paper makes no
claims about consciousness, sentience, or self-modeling beyond
what is already required to maintain agent identity across
sessions. The brain analogies are mapped to engineering
functions and are not load-bearing for the proposal.

\paragraph{Edge as a positioning choice.} We argue that the
contribution of an edge implementation is not capability
saturation but \emph{replicability}. The runtime described here
is intended to be reproducible by a single researcher with
hardware that costs less than two hundred dollars. Cluster
implementations of related ideas already exist in well-funded
laboratories and do not always ship.

\section{Conclusion}

The bottleneck in current multi-agent language model systems is
the convention that all inter-module communication must
serialize through human language. This convention is a holdover
from the chatbot era, where the only consumer was a human. When
the consumer is another module, the convention is wasteful and
lossy. We propose a runtime architecture, the thalamus layer,
that removes text from the inter-module bus by routing
modality-preserving vectors through a software router with a
layered hippocampal memory and an idle-time replay loop. We
describe a Phase 1 reference implementation deployable on a
Raspberry Pi 5 and a roadmap toward distributed and
hardware-accelerated versions. The thesis we defend is not that
the brain analogies are mechanically literal; it is that the
shape of the brain's modular processing offers an engineering
template that is well-matched to where multi-agent runtimes are
already heading, and that a small, open, replicable
implementation is a more useful contribution at this point than
another paper-only proposal.

\section*{Code and Reproducibility}

The reference implementation is available at:
\begin{itemize}[leftmargin=2em,itemsep=0pt]
  \item \url{https://github.com/msbel5/openclaw-thalamus} (this work, MIT)
  \item \url{https://github.com/msbel5/openclaw-aegis-signer} (companion: tool-call audit)
  \item \url{https://github.com/msbel5/openclaw-sga-mcts-atoms} (companion: atom retrieval)
\end{itemize}

\noindent
The corresponding npm packages are
\texttt{@msbel/openclaw-thalamus},
\texttt{@msbel/openclaw-aegis-signer}, and
\texttt{@msbel/openclaw-sga-mcts-atoms}, all released under
the MIT license. The longer-form, less formal version of this
proposal is available at \url{https://msbel.com}.

\section*{Acknowledgments}

The author thanks the maintainers of the OpenClaw project for
the plugin SDK that made this work possible. The author also
thanks Claude Sonnet 4.6 (Anthropic) and GPT-5 (OpenAI) for
sharpening the architectural framing during early drafts;
errors of judgment that survived this process remain the
author's.

\begin{thebibliography}{99}

\bibitem{aegis}
Bel, M.~S. (2026).
\emph{openclaw-aegis-signer: Ed25519-signed tool-call audit for OpenClaw}.
\url{https://github.com/msbel5/openclaw-aegis-signer}.

\bibitem{sgamcts}
Bel, M.~S. (2026).
\emph{openclaw-sga-mcts-atoms: Plan-time atom retrieval for OpenClaw}.
\url{https://github.com/msbel5/openclaw-sga-mcts-atoms}.

\bibitem{memlance}
OpenClaw bundled plugin. \emph{memory-lancedb}.
\url{https://docs.openclaw.ai/plugins/memory-lancedb}.

\bibitem{memwiki}
OpenClaw bundled plugin. \emph{memory-wiki}.
\url{https://docs.openclaw.ai/plugins/memory-wiki}.

\bibitem{siglip}
Zhai, X., Mustafa, B., Kolesnikov, A., \& Beyer, L. (2023).
Sigmoid Loss for Language Image Pre-Training (SigLIP).

\bibitem{dinov2}
Oquab, M., Darcet, T., Moutakanni, T., et~al. (2024).
DINOv2: Learning Robust Visual Features without Supervision.

\bibitem{whisper}
Radford, A., Kim, J.~W., Xu, T., et~al. (2023).
Robust Speech Recognition via Large-Scale Weak Supervision.

\bibitem{encodec}
Défossez, A., Copet, J., Synnaeve, G., \& Adi, Y. (2023).
High Fidelity Neural Audio Compression.

\bibitem{e5}
Wang, L., Yang, N., Huang, X., et~al. (2024).
Multilingual E5 Text Embeddings.

\bibitem{hipporag}
Gutiérrez, B.~J., Shu, Y., Gu, Y., Yasunaga, M., \& Su, Y. (2024).
HippoRAG: Neurobiologically Inspired Long-Term Memory for Large
Language Models.
\emph{NeurIPS}.
arXiv:2405.14831.

\bibitem{theaterofmind}
``Theater of Mind'' for LLMs (2026).
A Cognitive Architecture Based on Global Workspace Theory.
arXiv:2604.08206.

\bibitem{latentmas}
Latent Collaboration in Multi-Agent Systems (2025).
arXiv:2511.20639.

\bibitem{vjepa}
Bardes, A., Garrido, Q., Ponce, J., et~al. (2024).
V-JEPA: Latent Video Prediction for Visual Representation Learning.
\emph{OpenReview} WFYbBOEOtv.

\bibitem{perceiverio}
Jaegle, A., Borgeaud, S., Alayrac, J.-B., et~al. (2021).
Perceiver IO: A General Architecture for Structured Inputs and Outputs.
arXiv:2107.14795.

\bibitem{jepa}
LeCun, Y. (2022).
A Path Towards Autonomous Machine Intelligence.
\emph{OpenReview}.

\bibitem{dmn-review}
Raichle, M.~E. (2015).
The Brain's Default Mode Network.
\emph{Annual Review of Neuroscience}, 38, 433--447.

\end{thebibliography}

\end{document}
